% Fiches de lecture générées

\section{Développer la pratique réflexive dans le métier d'enseignant (Perrenoud, 2001)}
\textbf{Clé :} \texttt{perrenoud\_pratique\_reflexive}
\begin{itemize}
    \item \textbf{Mots-clés :} Pratique réflexive, professionnalisation, analyse de pratiques, savoirs d'action.
    \item \textbf{Résumé :} Philippe Perrenoud théorise l'enseignant comme un "praticien réflexif", capable d'analyser son action pendant (réflexion dans l'action) et après (réflexion sur l'action) pour s'adapter à des situations complexes et singulières. Il s'oppose à l'application mécanique de recettes pour prôner une intelligence de la situation.
    \item \textbf{Utilité VAE :} Fondamental pour justifier votre posture. Votre dossier ne se contente pas de décrire ce que vous avez fait, mais \textit{analyse} pourquoi et comment, montrant votre capacité à apprendre de votre expérience (notamment dans le Bilan et la Conclusion).
    \item \textbf{Citation :} \textit{"La pratique réflexive est un rapport au monde, une posture, une façon d'habiter son métier."}
\end{itemize}
\hrule

\section{La transposition didactique (Chevallard, 1985)}
\textbf{Clé :} \texttt{chevallard\_transposition}
\begin{itemize}
    \item \textbf{Mots-clés :} Savoir savant, savoir enseigné, noosphère, vigilance épistémologique.
    \item \textbf{Résumé :} Yves Chevallard décrit le processus par lequel un "savoir savant" (universitaire, scientifique) est transformé pour devenir un "savoir enseigné" (scolaire). Cette transformation est nécessaire mais comporte des risques de dénaturation (vigilance épistémologique).
    \item \textbf{Utilité VAE :} Essentiel pour l'Atelier IA / \LaTeX{} (Chap 9). Vous montrez comment vous avez transformé les concepts abstraits de la syntaxe \LaTeX{} et des LLMs en un dispositif concret d'écriture assistée par IA, opérant une transposition didactique interne consciente.
    \item \textbf{Citation :} \textit{"Un savoir n'est pas un objet qui se transporte d'un lieu à un autre sans subir de modifications."}
\end{itemize}
\hrule

\section{Les gestes professionnels et le jeu des postures (Bucheton \& Soulé, 2009)}
\textbf{Clé :} \texttt{bucheton\_multi\_agenda}
\begin{itemize}
    \item \textbf{Mots-clés :} Multi-agenda, étayage, tissage, atmosphère, pilotage.
    \item \textbf{Résumé :} Dominique Bucheton modélise l'activité enseignante comme un "multi-agenda" de préoccupations enchâssées (pilotage, atmosphère, tissage, étayage, objets de savoir). L'enseignant doit ajuster ses postures (contrôle, accompagnement, lâcher-prise...) en temps réel selon les réactions des élèves.
    \item \textbf{Utilité VAE :} Utilisé explicitement dans le Chapitre 10 (Jeu Vidéo). Vous justifiez votre gestion de classe complexe en expliquant comment vous avez jonglé entre "faire avancer le projet" (pilotage) et "soutenir les élèves en difficulté" (étayage).
    \item \textbf{Citation :} \textit{"Enseigner, c'est ajuster."}
\end{itemize}
\hrule

\section{Mind in Society: The Development of Higher Psychological Processes (Vygotsky, 1978)}
\textbf{Clé :} \texttt{vygotsky\_mind\_society}
\begin{itemize}
    \item \textbf{Mots-clés :} Zone Proximale de Développement (ZPD), médiation, socioconstructivisme, langage intérieur.
    \item \textbf{Résumé :} Lev Vygotsky pose les bases du socioconstructivisme : l'apprentissage est d'abord social avant d'être individuel. Il introduit la ZPD, l'écart entre ce que l'élève peut faire seul et ce qu'il peut faire avec aide. L'enseignant agit comme médiateur dans cette zone.
    \item \textbf{Utilité VAE :} Justifie vos travaux de groupe (Jeu Vidéo) et l'usage de l'IA (Thèse). L'IA peut être vue comme un outil de médiation permettant à l'élève d'atteindre sa ZPD.
    \item \textbf{Citation :} \textit{"Ce que l'enfant sait faire aujourd'hui en collaboration, il saura le faire tout seul demain."}
\end{itemize}
\hrule

\section{The Reflective Practitioner (Schön, 1983)}
\textbf{Clé :} \texttt{schon\_reflective\_practitioner}
\begin{itemize}
    \item \textbf{Mots-clés :} Réflexion dans l'action, réflexion sur l'action, savoir tacite.
    \item \textbf{Résumé :} Donald Schön étudie comment les professionnels (architectes, médecins...) agissent en situation d'incertitude. Ils ne se contentent pas d'appliquer des théories, mais conversent avec la situation.
    \item \textbf{Utilité VAE :} Renforce la référence à Perrenoud. Valide votre expérience industrielle comme source de savoirs : vous avez développé une expertise par la pratique que vous formalisez maintenant.
    \item \textbf{Citation :} \textit{"Le praticien réflexif converse avec la situation."}
\end{itemize}
\hrule

\section{Experience and Education (Dewey, 1938)}
\textbf{Clé :} \texttt{dewey\_experience\_education}
\begin{itemize}
    \item \textbf{Mots-clés :} Learning by doing, expérience, continuité, interaction.
    \item \textbf{Résumé :} John Dewey critique l'école traditionnelle et prône une éducation fondée sur l'expérience vécue par l'élève. L'apprentissage doit avoir du sens et être connecté à la réalité sociale.
    \item \textbf{Utilité VAE :} Justifie votre approche par projet (Jeu Vidéo) et par investigation (SQL). Vous privilégiez le "faire" pour comprendre.
    \item \textbf{Citation :} \textit{"Toute éducation authentique provient de l'expérience."}
\end{itemize}
\hrule

\section{Constructivisme et socioconstructivisme (Jonnaert, 2002)}
\textbf{Clé :} \texttt{constructivisme\_education}
\begin{itemize}
    \item \textbf{Mots-clés :} Construction des savoirs, conflit sociocognitif, représentations.
    \item \textbf{Résumé :} Synthèse des théories constructivistes (Piaget) et socioconstructivistes. L'élève construit ses connaissances en remettant en cause ses représentations initiales (conflit cognitif) et en interagissant avec les autres.
    \item \textbf{Utilité VAE :} Cadre théorique global de votre pédagogie. Vous ne "transmettez" pas le savoir, vous créez les conditions pour que les élèves le construisent (ex: rétro-ingénierie).
\end{itemize}
\hrule

\section{Artificial Intelligence in Education: Promises and Implications (Luckin et al., 2022)}
\textbf{Clé :} \texttt{luckin\_ai\_education}
\begin{itemize}
    \item \textbf{Mots-clés :} IA, personnalisation, tuteur intelligent, éthique, rôle de l'enseignant.
    \item \textbf{Résumé :} Rosemary Luckin, experte mondiale de l'IA en éducation, analyse comment l'IA peut soutenir (et non remplacer) l'enseignant. Elle distingue trois types d'IA : pour l'apprenant (tuteur), pour l'enseignant (assistant) et pour le système (gestion). Elle insiste sur l'importance de l'explicabilité (XAI).
    \item \textbf{Utilité VAE :} Cœur théorique de votre projet de thèse (Chap 17). Vous utilisez ses concepts pour justifier que l'IA doit être un "assistant" pour l'enseignant (pour l'adaptation des ressources) et non une "boîte noire".
    \item \textbf{Citation :} \textit{"AI should be designed to empower teachers, not replace them."}
\end{itemize}
\hrule

\section{AI and Education: Guidance for Policy Makers (UNESCO, 2019)}
\textbf{Clé :} \texttt{unesco\_ai\_education}
\begin{itemize}
    \item \textbf{Mots-clés :} Politiques publiques, équité, inclusion, formation des enseignants.
    \item \textbf{Résumé :} Rapport de référence de l'UNESCO qui pose les cadres éthiques et politiques de l'IA. Il met en garde contre les risques de fracture numérique et d'amplification des biais. Il prône une approche "human-centered".
    \item \textbf{Utilité VAE :} Donne une dimension institutionnelle et internationale à votre réflexion. Vous montrez que votre projet s'inscrit dans les grandes orientations mondiales (ODD 4).
\end{itemize}
\hrule

\section{Application and theory gaps during the rise of AI in education (Chen et al., 2022)}
\textbf{Clé :} \texttt{chen\_ai\_personalized\_learning}
\begin{itemize}
    \item \textbf{Mots-clés :} Revue de littérature, fossé théorique, technologies émergentes.
    \item \textbf{Résumé :} Cette méta-analyse montre qu'il y a souvent un décalage entre les avancées technologiques de l'IA (Deep Learning) et leur application pédagogique réelle. Beaucoup d'outils sont techniquement impressionnants mais pédagogiquement pauvres.
    \item \textbf{Utilité VAE :} Justifie votre posture de "praticien-chercheur" : vous êtes là pour combler ce fossé, en apportant votre expertise terrain à la conception d'outils IA.
\end{itemize}
\hrule

\section{Intelligence artificielle et éducation (Académie des sciences, 2023)}
\textbf{Clé :} \texttt{intelligence\_artificielle\_education}
\begin{itemize}
    \item \textbf{Mots-clés :} Souveraineté numérique, formation, pensée critique.
    \item \textbf{Résumé :} Avis de l'Académie des sciences française. Insiste sur la nécessité de former les élèves \textit{à} l'IA (comprendre comment ça marche) et pas seulement \textit{par} l'IA.
    \item \textbf{Utilité VAE :} Valide votre approche en NSI : vous enseignez le fonctionnement des algorithmes (Chap 10, Jeu Vidéo, IA des ennemis) pour démystifier l'outil.
\end{itemize}
\hrule

\section{Éduquer au numérique (Henry \& Boraita, 2021)}
\textbf{Clé :} \texttt{henry\_bora\_2021}
\begin{itemize}
    \item \textbf{Mots-clés :} Culture numérique, citoyenneté, pensée informatique.
    \item \textbf{Résumé :} Ouvrage de vulgarisation et de didactique qui donne des clés pour enseigner l'informatique sans se noyer dans la technique. Met l'accent sur les enjeux sociétaux.
    \item \textbf{Utilité VAE :} Utile pour la partie SNT (Sciences Numériques et Technologie) où l'enjeu est la culture générale numérique plus que le code pur.
\end{itemize}
\hrule

\section{Fiche RNCP38356 - Master MEEF Informatique (France Compétences, 2024)}
\textbf{Clé :} \texttt{rncp38356}
\begin{itemize}
    \item \textbf{Mots-clés :} Compétences, certification, cadre légal.
    \item \textbf{Résumé :} Document officiel définissant les compétences attendues pour le Master MEEF. Il structure votre dossier VAE (les blocs de compétences).
    \item \textbf{Utilité VAE :} C'est la "bible" de votre VAE. Vous devez montrer que vous cochez toutes les cases de ce référentiel.
\end{itemize}
\hrule

\section{Référentiel des compétences professionnelles des métiers du professorat (MEN, 2013)}
\textbf{Clé :} \texttt{referentiel\_competences}
\begin{itemize}
    \item \textbf{Mots-clés :} 19 compétences, éthique, langue, numérique.
    \item \textbf{Résumé :} Texte fondateur définissant ce qu'est un enseignant en France. Il liste les compétences communes (CC) et spécifiques (P).
    \item \textbf{Utilité VAE :} Vous y faites référence pour prouver que vous avez la posture d'un fonctionnaire d'État (CC1 : Faire partager les valeurs de la République).
\end{itemize}
\hrule

\section{Programmes NSI et SNT (Éduscol, 2019/2024)}
\textbf{Clé :} \texttt{eduscol\_nsi}
\begin{itemize}
    \item \textbf{Mots-clés :} Curriculum, attendus, progression.
    \item \textbf{Résumé :} Les programmes officiels que vous enseignez. Ils définissent les contenus (données, algo, langages, machines...).
    \item \textbf{Utilité VAE :} Prouve votre maîtrise disciplinaire. Vous connaissez le "texte du savoir" que vous devez enseigner.
\end{itemize}
\hrule

\section{La société inclusive, parlons-en ! (Gardou, 2012)}
\textbf{Clé :} \texttt{gardou\_inclusive\_society}
\begin{itemize}
    \item \textbf{Mots-clés :} Inclusion, handicap, accessibilité universelle.
    \item \textbf{Résumé :} Charles Gardou défend une société où l'adaptation n'est pas l'effort du handicapé, mais celui de la société. Il ne s'agit pas d'intégrer (mettre dedans) mais d'inclure (adapter l'environnement).
    \item \textbf{Utilité VAE :} Justifie vos efforts pour l'accessibilité numérique (Chap 6, Forge) et la différenciation pédagogique (Chap 10).
\end{itemize}
\hrule

\section{Différenciation pédagogique (CNESCO, 2017)}
\textbf{Clé :} \texttt{cnesco\_inclusion}
\begin{itemize}
    \item \textbf{Mots-clés :} Hétérogénéité, flexibilité, réussite pour tous.
    \item \textbf{Résumé :} Dossier de synthèse sur comment gérer l'hétérogénéité des élèves. Prône la variation des supports, des processus et des productions.
    \item \textbf{Utilité VAE :} Valide votre approche dans le projet Jeu Vidéo (groupes de niveaux, rôles différents) pour gérer une classe hétérogène.
\end{itemize}
\hrule

\section{Documentation Python / SQL / W3C (Tech Refs)}
\textbf{Clé :} \texttt{python\_education}, \texttt{sql\_education}, \texttt{accessibilite\_numerique}
\begin{itemize}
    \item \textbf{Mots-clés :} Standards, documentation technique, normes.
    \item \textbf{Résumé :} Documentations officielles des technologies enseignées.
    \item \textbf{Utilité VAE :} Montre que vous vous formez aux sources primaires et que vous enseignez des standards industriels reconnus, pas des "bricolages".
\end{itemize}
\hrule
